% Options for packages loaded elsewhere
% Options for packages loaded elsewhere
\PassOptionsToPackage{unicode}{hyperref}
\PassOptionsToPackage{hyphens}{url}
\PassOptionsToPackage{dvipsnames,svgnames,x11names}{xcolor}
%
\documentclass[
  a4paper,
  DIV=11,
  numbers=noendperiod,
  oneside]{scrreport}
\usepackage{xcolor}
\usepackage[left=10mm,marginparwidth=60mm,textwidth=130mm,marginparsep=5mm]{geometry}
\usepackage{amsmath,amssymb}
\setcounter{secnumdepth}{-\maxdimen} % remove section numbering
\usepackage{iftex}
\ifPDFTeX
  \usepackage[T1]{fontenc}
  \usepackage[utf8]{inputenc}
  \usepackage{textcomp} % provide euro and other symbols
\else % if luatex or xetex
  \usepackage{unicode-math} % this also loads fontspec
  \defaultfontfeatures{Scale=MatchLowercase}
  \defaultfontfeatures[\rmfamily]{Ligatures=TeX,Scale=1}
\fi
\usepackage{lmodern}
\ifPDFTeX\else
  % xetex/luatex font selection
\fi
% Use upquote if available, for straight quotes in verbatim environments
\IfFileExists{upquote.sty}{\usepackage{upquote}}{}
\IfFileExists{microtype.sty}{% use microtype if available
  \usepackage[]{microtype}
  \UseMicrotypeSet[protrusion]{basicmath} % disable protrusion for tt fonts
}{}
\makeatletter
\@ifundefined{KOMAClassName}{% if non-KOMA class
  \IfFileExists{parskip.sty}{%
    \usepackage{parskip}
  }{% else
    \setlength{\parindent}{0pt}
    \setlength{\parskip}{6pt plus 2pt minus 1pt}}
}{% if KOMA class
  \KOMAoptions{parskip=half}}
\makeatother
% Make \paragraph and \subparagraph free-standing
\makeatletter
\ifx\paragraph\undefined\else
  \let\oldparagraph\paragraph
  \renewcommand{\paragraph}{
    \@ifstar
      \xxxParagraphStar
      \xxxParagraphNoStar
  }
  \newcommand{\xxxParagraphStar}[1]{\oldparagraph*{#1}\mbox{}}
  \newcommand{\xxxParagraphNoStar}[1]{\oldparagraph{#1}\mbox{}}
\fi
\ifx\subparagraph\undefined\else
  \let\oldsubparagraph\subparagraph
  \renewcommand{\subparagraph}{
    \@ifstar
      \xxxSubParagraphStar
      \xxxSubParagraphNoStar
  }
  \newcommand{\xxxSubParagraphStar}[1]{\oldsubparagraph*{#1}\mbox{}}
  \newcommand{\xxxSubParagraphNoStar}[1]{\oldsubparagraph{#1}\mbox{}}
\fi
\makeatother

\usepackage{color}
\usepackage{fancyvrb}
\newcommand{\VerbBar}{|}
\newcommand{\VERB}{\Verb[commandchars=\\\{\}]}
\DefineVerbatimEnvironment{Highlighting}{Verbatim}{commandchars=\\\{\}}
% Add ',fontsize=\small' for more characters per line
\usepackage{framed}
\definecolor{shadecolor}{RGB}{241,243,245}
\newenvironment{Shaded}{\begin{snugshade}}{\end{snugshade}}
\newcommand{\AlertTok}[1]{\textcolor[rgb]{0.68,0.00,0.00}{#1}}
\newcommand{\AnnotationTok}[1]{\textcolor[rgb]{0.37,0.37,0.37}{#1}}
\newcommand{\AttributeTok}[1]{\textcolor[rgb]{0.40,0.45,0.13}{#1}}
\newcommand{\BaseNTok}[1]{\textcolor[rgb]{0.68,0.00,0.00}{#1}}
\newcommand{\BuiltInTok}[1]{\textcolor[rgb]{0.00,0.23,0.31}{#1}}
\newcommand{\CharTok}[1]{\textcolor[rgb]{0.13,0.47,0.30}{#1}}
\newcommand{\CommentTok}[1]{\textcolor[rgb]{0.37,0.37,0.37}{#1}}
\newcommand{\CommentVarTok}[1]{\textcolor[rgb]{0.37,0.37,0.37}{\textit{#1}}}
\newcommand{\ConstantTok}[1]{\textcolor[rgb]{0.56,0.35,0.01}{#1}}
\newcommand{\ControlFlowTok}[1]{\textcolor[rgb]{0.00,0.23,0.31}{\textbf{#1}}}
\newcommand{\DataTypeTok}[1]{\textcolor[rgb]{0.68,0.00,0.00}{#1}}
\newcommand{\DecValTok}[1]{\textcolor[rgb]{0.68,0.00,0.00}{#1}}
\newcommand{\DocumentationTok}[1]{\textcolor[rgb]{0.37,0.37,0.37}{\textit{#1}}}
\newcommand{\ErrorTok}[1]{\textcolor[rgb]{0.68,0.00,0.00}{#1}}
\newcommand{\ExtensionTok}[1]{\textcolor[rgb]{0.00,0.23,0.31}{#1}}
\newcommand{\FloatTok}[1]{\textcolor[rgb]{0.68,0.00,0.00}{#1}}
\newcommand{\FunctionTok}[1]{\textcolor[rgb]{0.28,0.35,0.67}{#1}}
\newcommand{\ImportTok}[1]{\textcolor[rgb]{0.00,0.46,0.62}{#1}}
\newcommand{\InformationTok}[1]{\textcolor[rgb]{0.37,0.37,0.37}{#1}}
\newcommand{\KeywordTok}[1]{\textcolor[rgb]{0.00,0.23,0.31}{\textbf{#1}}}
\newcommand{\NormalTok}[1]{\textcolor[rgb]{0.00,0.23,0.31}{#1}}
\newcommand{\OperatorTok}[1]{\textcolor[rgb]{0.37,0.37,0.37}{#1}}
\newcommand{\OtherTok}[1]{\textcolor[rgb]{0.00,0.23,0.31}{#1}}
\newcommand{\PreprocessorTok}[1]{\textcolor[rgb]{0.68,0.00,0.00}{#1}}
\newcommand{\RegionMarkerTok}[1]{\textcolor[rgb]{0.00,0.23,0.31}{#1}}
\newcommand{\SpecialCharTok}[1]{\textcolor[rgb]{0.37,0.37,0.37}{#1}}
\newcommand{\SpecialStringTok}[1]{\textcolor[rgb]{0.13,0.47,0.30}{#1}}
\newcommand{\StringTok}[1]{\textcolor[rgb]{0.13,0.47,0.30}{#1}}
\newcommand{\VariableTok}[1]{\textcolor[rgb]{0.07,0.07,0.07}{#1}}
\newcommand{\VerbatimStringTok}[1]{\textcolor[rgb]{0.13,0.47,0.30}{#1}}
\newcommand{\WarningTok}[1]{\textcolor[rgb]{0.37,0.37,0.37}{\textit{#1}}}

\usepackage{longtable,booktabs,array}
\usepackage{calc} % for calculating minipage widths
% Correct order of tables after \paragraph or \subparagraph
\usepackage{etoolbox}
\makeatletter
\patchcmd\longtable{\par}{\if@noskipsec\mbox{}\fi\par}{}{}
\makeatother
% Allow footnotes in longtable head/foot
\IfFileExists{footnotehyper.sty}{\usepackage{footnotehyper}}{\usepackage{footnote}}
\makesavenoteenv{longtable}
\usepackage{graphicx}
\makeatletter
\newsavebox\pandoc@box
\newcommand*\pandocbounded[1]{% scales image to fit in text height/width
  \sbox\pandoc@box{#1}%
  \Gscale@div\@tempa{\textheight}{\dimexpr\ht\pandoc@box+\dp\pandoc@box\relax}%
  \Gscale@div\@tempb{\linewidth}{\wd\pandoc@box}%
  \ifdim\@tempb\p@<\@tempa\p@\let\@tempa\@tempb\fi% select the smaller of both
  \ifdim\@tempa\p@<\p@\scalebox{\@tempa}{\usebox\pandoc@box}%
  \else\usebox{\pandoc@box}%
  \fi%
}
% Set default figure placement to htbp
\def\fps@figure{htbp}
\makeatother





\setlength{\emergencystretch}{3em} % prevent overfull lines

\providecommand{\tightlist}{%
  \setlength{\itemsep}{0pt}\setlength{\parskip}{0pt}}



 


% This is used for pdf compilation from qmd.
% Any special LaTeX commands should be updated in parallel with mymacros.sty which is used for html compilation.
% All rows *might* need an HTML equivalent. Those at the bottom of this file are identical.
% This is how the two files are included in the _quarto.yaml header of the qmd file
%format:
%  html:
%    defaults: [resources/latex/global-maths-definitions.yaml]
%  pdf:
%   include-in-header: ["include-in-header.tex"]

% the 'defaults' above secretly points to mymacros.sty where they really live.

\usepackage{gensymb}
\usepackage{amsmath}
\usepackage{xcolor}

\usepackage[skins,breakable]{tcolorbox}

% Code to reduce width of code blocks in PDF to encourage line wrapping
\usepackage{fvextra}
      \DefineVerbatimEnvironment{Highlighting}{Verbatim}{
        breaklines,
        breakanywhere,
        commandchars=\\\{\},
        fontsize=\small
      }
% Also line wrap in verbatim blocks
\RecustomVerbatimEnvironment{verbatim}{Verbatim}{breaklines,fontsize=\footnotesize}

\usepackage{environ}

% Below here is an old way to defining new custom coloured boxes.
% We use nice extensions now.

% \definecolor{goals-col}{RGB}{201, 232, 222}
% 
% % For defining new callout boxes
% % HTML version is in CSS file
% % Need to turn on goals and goals-header as environments using the latex-environment extension for quarto
% % This is designed to work for the following fenced div syntax:
% % :::{.goals}
% % ::::{.goals-header}
% % TITLE
% % ::::
% % Actual box contents
% % :::
% \newtcolorbox{goals}[1][]{notitle, top=0pt, boxsep=0pt, coltitle=black!90!white, colback=white, colframe=goals-col,#1}
% \NewEnviron{goals-header}{\tcbsubtitle[top=1mm, bottom=1mm]{\BODY}}
% % It works by creating a box with no title. Then no spacing. Then a subtitle defined later.

% Increasing spacing between rows of tables
\def\arraystretch{2}

% Syntax for highlighting my KEY TERMS
% HTML version is included in css for .term
\newcommand{\term}[1]{\textbf{#1}}

% Added to do derivatives, custom commands for HTML
% read by MathJax are inserted via diff.macro
% must keep naming in sync with esdiff package
\usepackage{esdiff}

%%%%%%%%%%%%%%%%%%%%%%%%%%%%%%%%%%%%%%%%%%%%%%%%%%%%%%%%%%%%
%%% BELOW HERE ARE COMMANDS SHARED WITH THE HTML OUTPUT %%%%
%%%%%%%%%%%%%%%%%%%%%%%%%%%%%%%%%%%%%%%%%%%%%%%%%%%%%%%%%%%%
%%%% KEEP THESE UPDATED ALSO IN THE MYMACROS.STY FILE   %%%%
%%%%%%%%%%%%%%%%%%%%%%%%%%%%%%%%%%%%%%%%%%%%%%%%%%%%%%%%%%%%

% Adding some operator names missing by default
\DeclareMathOperator{\cosec}{cosec}
\DeclareMathOperator{\powop}{pow} %\newcommand{\pow}{\operatorname{pow}}
\definecolor{frog}{rgb}{1,0,0}
\newcommand{\dx}{\,\text{d}x}
\newcommand{\du}{\,\text{d}u}
\newcommand{\dy}{\,\text{d}y}
\newcommand{\dt}{\,\text{d}t}

\KOMAoption{captions}{tableheading}
\makeatletter
\@ifpackageloaded{float}{}{\usepackage{float}}
\floatstyle{plain}
\@ifundefined{c@chapter}{\newfloat{video}{h}{lovid}}{\newfloat{video}{h}{lovid}[chapter]}
\floatname{video}{Video}
\newcommand*\listofvideos{\listof{video}{List of Videos}}
\makeatother
\makeatletter
\@ifpackageloaded{tcolorbox}{}{\usepackage[skins,breakable]{tcolorbox}}
\@ifpackageloaded{fontawesome5}{}{\usepackage{fontawesome5}}
\definecolor{quarto-callout-color}{HTML}{909090}
\definecolor{quarto-callout-note-color}{HTML}{0758E5}
\definecolor{quarto-callout-important-color}{HTML}{CC1914}
\definecolor{quarto-callout-warning-color}{HTML}{EB9113}
\definecolor{quarto-callout-tip-color}{HTML}{00A047}
\definecolor{quarto-callout-caution-color}{HTML}{FC5300}
\definecolor{quarto-callout-color-frame}{HTML}{acacac}
\definecolor{quarto-callout-note-color-frame}{HTML}{4582ec}
\definecolor{quarto-callout-important-color-frame}{HTML}{d9534f}
\definecolor{quarto-callout-warning-color-frame}{HTML}{f0ad4e}
\definecolor{quarto-callout-tip-color-frame}{HTML}{02b875}
\definecolor{quarto-callout-caution-color-frame}{HTML}{fd7e14}
\makeatother
\makeatletter
\@ifpackageloaded{caption}{}{\usepackage{caption}}
\AtBeginDocument{%
\ifdefined\contentsname
  \renewcommand*\contentsname{Table of contents}
\else
  \newcommand\contentsname{Table of contents}
\fi
\ifdefined\listfigurename
  \renewcommand*\listfigurename{List of Figures}
\else
  \newcommand\listfigurename{List of Figures}
\fi
\ifdefined\listtablename
  \renewcommand*\listtablename{List of Tables}
\else
  \newcommand\listtablename{List of Tables}
\fi
\ifdefined\figurename
  \renewcommand*\figurename{Figure}
\else
  \newcommand\figurename{Figure}
\fi
\ifdefined\tablename
  \renewcommand*\tablename{Table}
\else
  \newcommand\tablename{Table}
\fi
}
\@ifpackageloaded{float}{}{\usepackage{float}}
\floatstyle{ruled}
\@ifundefined{c@chapter}{\newfloat{codelisting}{h}{lop}}{\newfloat{codelisting}{h}{lop}[chapter]}
\floatname{codelisting}{Listing}
\newcommand*\listoflistings{\listof{codelisting}{List of Listings}}
\makeatother
\makeatletter
\makeatother
\makeatletter
\@ifpackageloaded{caption}{}{\usepackage{caption}}
\@ifpackageloaded{subcaption}{}{\usepackage{subcaption}}
\makeatother
\makeatletter
\@ifpackageloaded{sidenotes}{}{\usepackage{sidenotes}}
\@ifpackageloaded{marginnote}{}{\usepackage{marginnote}}
\makeatother
\makeatletter
\@ifpackageloaded{tcolorbox}{}{\usepackage[many]{tcolorbox}}
\makeatother
%%%% ---foldboxy preamble ----- %%%%%

\definecolor{fbx-default-color1}{HTML}{c7c7d0}
\definecolor{fbx-default-color2}{HTML}{a3a3aa}

\definecolor{fbox-color1}{HTML}{c7c7d0}
\definecolor{fbox-color2}{HTML}{a3a3aa}

% arguments: #1 typelabelnummer: #2 titel: #3
\newenvironment{fbx}[3]{\begin{tcolorbox}[enhanced, breakable,%
attach boxed title to top*={xshift=1.4pt},
boxed title style={boxrule=0.0mm, fuzzy shadow={1pt}{-1pt}{0mm}{0.1mm}{gray}, arc=.3em, rounded corners=east, sharp corners=west}, colframe=#1-color2, colbacktitle=#1-color1, colback = white, coltitle=black,  titlerule=0mm, toprule=0pt, bottomrule=.7pt, leftrule=.3em, rightrule=0pt, outer arc=.3em,  arc=0pt,	 sharp corners = east, left=.5em, bottomtitle=1mm, toptitle=1mm,title=\textbf{#2}\hspace{0.5em}{#3}]}
{\end{tcolorbox}}

% boxed environment with right border
\newenvironment{fbxSimple}[3]{\begin{tcolorbox}[enhanced, breakable,%
attach boxed title to top*={xshift=1.4pt},
boxed title style={boxrule=0.0mm, fuzzy shadow={1pt}{-1pt}{0mm}{0.1mm}{gray}, arc=.3em, rounded corners=east, sharp corners=west}, colframe=#1-color2, colbacktitle=#1-color1, colback = white, coltitle=black,  titlerule=0mm, toprule=0pt, bottomrule=.7pt, leftrule=.3em, rightrule=.7pt, outer arc=.3em,  	left=.5em, right=.5em, bottomtitle=1mm, toptitle=1mm,title=\textbf{#2}\hspace{0.5em}{#3}]}
{\end{tcolorbox}}

%%%% --- end foldboxy preamble ----- %%%%%
%%==== colors from yaml ===%
\definecolor{Corollary-color1}{HTML}{7fbad7}
\definecolor{Corollary-color2}{HTML}{005398}
\definecolor{Weblink-color1}{HTML}{d18888}
\definecolor{Weblink-color2}{HTML}{660000}
\definecolor{Lemma-color1}{HTML}{b2dac4}
\definecolor{Lemma-color2}{HTML}{004721}
\definecolor{Video-color1}{HTML}{d18888}
\definecolor{Video-color2}{HTML}{660000}
\definecolor{Answer-color1}{HTML}{d18888}
\definecolor{Answer-color2}{HTML}{660000}
\definecolor{Example-color1}{HTML}{b2dac4}
\definecolor{Example-color2}{HTML}{004721}
\definecolor{Definition-color1}{HTML}{ffdc36}
\definecolor{Definition-color2}{HTML}{ffb948}
\definecolor{Task-color1}{HTML}{948bde}
\definecolor{Task-color2}{HTML}{584eab}
\definecolor{Theorem-color1}{HTML}{948bde}
\definecolor{Theorem-color2}{HTML}{584eab}
%=============%
\usepackage{bookmark}
\IfFileExists{xurl.sty}{\usepackage{xurl}}{} % add URL line breaks if available
\urlstyle{same}
\hypersetup{
  colorlinks=true,
  linkcolor={blue},
  filecolor={Maroon},
  citecolor={Blue},
  urlcolor={Blue},
  pdfcreator={LaTeX via pandoc}}


\author{}
\date{}
\begin{document}


\chapter{Chapter 4 - Nicer boxes}\label{chapter-4---nicer-boxes}

This chapter shows off some nice coloured boxes. They come from the
\texttt{custom-numbered-blocks} extension. First you need to install the
extension \texttt{quarto\ add\ ute/custom-numbered-blocks}.

This code has then been added to the \texttt{\_quarto.yml} or
\texttt{\_metadata.yml} file (or at the top of this file) to load the
coloured boxes.

\begin{verbatim}
filters:
  - custom-numbered-blocks # This loads the extension

# This is ute-hahn's nice coloured boxes: setup
custom-numbered-blocks:
  groups:
    thmlike:
      colors: [948bde, 584eab]
      boxstyle: foldbox.simple
    deflike:
      colors: [ffdc36, ffb948]
      boxstyle: foldbox.simple
    videolike:
      colors: [beb9eb, 877eda]
      boxstyle: foldbox.simple
      collapse: open
  classes:
    Theorem:
      group: thmlike
    Lemma:
      group: thmlike
    Corollary:
      group: thmlike
    Definition:
      group: deflike
    Video:
      group: videolike
      
crossref:
  chapters: true
\end{verbatim}

\section{Custom Numbered Blocks}\label{custom-numbered-blocks}

We've defined several custom-numbered block types using the
\texttt{custom-numbered-blocks} option. The key point is that these can
be numbered and follow grouped numbering conventions.

Each block belongs to a \textbf{group} that determines its shared
numbering.

All of these blocks are \emph{foldable}, meaning they can be opened and
closed. The \texttt{collapse} setting determines the starting state.

\begin{longtable}[]{@{}
  >{\raggedright\arraybackslash}p{(\linewidth - 10\tabcolsep) * \real{0.1287}}
  >{\raggedright\arraybackslash}p{(\linewidth - 10\tabcolsep) * \real{0.1188}}
  >{\raggedright\arraybackslash}p{(\linewidth - 10\tabcolsep) * \real{0.1881}}
  >{\raggedright\arraybackslash}p{(\linewidth - 10\tabcolsep) * \real{0.2376}}
  >{\raggedright\arraybackslash}p{(\linewidth - 10\tabcolsep) * \real{0.1683}}
  >{\raggedright\arraybackslash}p{(\linewidth - 10\tabcolsep) * \real{0.1584}}@{}}
\toprule\noalign{}
\begin{minipage}[b]{\linewidth}\raggedright
Box Name
\end{minipage} & \begin{minipage}[b]{\linewidth}\raggedright
Group
\end{minipage} & \begin{minipage}[b]{\linewidth}\raggedright
Left Border Color
\end{minipage} & \begin{minipage}[b]{\linewidth}\raggedright
Title Background Color
\end{minipage} & \begin{minipage}[b]{\linewidth}\raggedright
Default State
\end{minipage} & \begin{minipage}[b]{\linewidth}\raggedright
\end{minipage} \\
\midrule\noalign{}
\endhead
\bottomrule\noalign{}
\endlastfoot
Theorem & \texttt{thmlike} & \texttt{\#948bde} & \texttt{\#584eab} &
Open / Closeable & \\
Lemma & \texttt{thmlike} & \texttt{\#b2dac4} & \texttt{\#004721} & Open
/ Closeable & \\
Corollary & \texttt{thmlike} & \texttt{\#7fbad7} & \texttt{\#005398} &
Open / Closeable & \\
Definition & \texttt{deflike} & \texttt{\#ffdc36} & \texttt{\#ffb948} &
Open / Closeable & \\
Video & \texttt{videolike} & \texttt{\#d18888} & \texttt{\#660000} &
Open / Closeable & \\
\end{longtable}

\begin{quote}
Note: The two colors defined for each group represent the \textbf{left
border} (first color) and the \textbf{title background} (second color).
\end{quote}

These boxes also have dark-mode versions, see below.

\phantomsection\label{vidFirst}
\begin{fbxSimple}{Video}{Video 1: }{Binomial response models applied to beetle mortality data (9m59)}
\phantomsection\label{vidFirst}
There is some text here before the video

\url{https://youtu.be/d2nG7Y17sQw}

\end{fbxSimple}

\subsection{Beetle mortality}\label{beetle-mortality}

Data for this example consists of the number of beetles dead
(\emph{killed}) after five hours exposure to gaseous carbon disulphide
at various concentrations (\emph{dose}). The goal for this analysis is
to model the probability of a beetle dying as a function of the carbon
disulphide dose.

\subsection{R}

\begin{Shaded}
\begin{Highlighting}[]
\NormalTok{beetles }\OtherTok{\textless{}{-}} \FunctionTok{read.csv}\NormalTok{(}\FunctionTok{url}\NormalTok{(}\StringTok{"http://www.stats.gla.ac.uk/\textasciitilde{}tereza/rp/beetles.csv"}\NormalTok{))}
\NormalTok{beetles}
\end{Highlighting}
\end{Shaded}

\subsection{Python}

\begin{Shaded}
\begin{Highlighting}[]
\ImportTok{import}\NormalTok{ pandas }\ImportTok{as}\NormalTok{ pd}
\CommentTok{\# beetles = pd.read\_csv("resources/data/beetles.csv")}
\CommentTok{\# beetles}
\end{Highlighting}
\end{Shaded}

\begin{tcolorbox}[enhanced jigsaw, colbacktitle=quarto-callout-tip-color!10!white, rightrule=.15mm, toprule=.15mm, opacityback=0, title=\textcolor{quarto-callout-tip-color}{\faLightbulb}\hspace{0.5em}{Conventional tip callout}, colback=white, bottomrule=.15mm, colframe=quarto-callout-tip-color-frame, breakable, leftrule=.75mm, bottomtitle=1mm, opacitybacktitle=0.6, toptitle=1mm, left=2mm, titlerule=0mm, arc=.35mm, coltitle=black]

Some content

\end{tcolorbox}

\phantomsection\label{bestlemma}
\begin{fbxSimple}{Lemma}{Lemma 1: }{Super}
\phantomsection\label{bestlemma}
This is a fantastic Lemma.

\end{fbxSimple}

There are no Lemma's better than Lemma \hyperref[bestlemma]{1}.

The other way to embed a video is inside a fenced div which just
contains a label starting with \# like this:

Our first video was embedded inside a custom numbered block and it was
called Video \hyperref[vidFirst]{1}.

\section{Demonstration of Custom Numbered
Blocks}\label{demonstration-of-custom-numbered-blocks}

Below are examples of the custom-numbered blocks.\\
Each block is foldable (open by default) and has its own color scheme
and group style as defined in your YAML configuration.

Notice Theorem, Corollary and Lemma has been defined to share counter.
But Definition and Video have their own types (\emph{deflike} and
\emph{vidlike}).

\phantomsection\label{rthmpythag}
\begin{fbxSimple}{Theorem}{Theorem 2: }{Pythagoras’ Theorem}
\phantomsection\label{rthmpythag}
In a right-angled triangle\ldots{}

\end{fbxSimple}

\phantomsection\label{Rcor-pythag}
\begin{fbxSimple}{Corollary}{Corollary 3: }{Corollary to Pythagoras’ Theorem}
\phantomsection\label{Rcor-pythag}
In a right-angled triangle, if the two shorted sides are equal, so that
\(a = b\), then the hypotenuse satisfies \(c=a\sqrt{2}\).

\end{fbxSimple}

\phantomsection\label{Rlem-euclid}
\begin{fbxSimple}{Lemma}{Lemma 4: }{Euclid}
\phantomsection\label{Rlem-euclid}
If two lines are perpendicular to the same line, then they are parallel
to each other.

\end{fbxSimple}

\phantomsection\label{Rdef-orthog}
\begin{fbxSimple}{Definition}{Definition 1: }{Definition of Orthogonality}
\phantomsection\label{Rdef-orthog}
Two vectors \(\mathbf{u}\) and \(\mathbf{v}\) are said to be
\emph{orthogonal} if their dot product is zero: \[
\mathbf{u} \cdot \mathbf{v} = 0.
\]

\end{fbxSimple}

\phantomsection\label{Rvid-yt}
\begin{fbxSimple}{Video}{Video 2: }{Video Demonstration**}
\phantomsection\label{Rvid-yt}
Here's a Youtube video:

\url{https://youtu.be/d2nG7Y17sQw}

\end{fbxSimple}

\begin{center}\rule{0.5\linewidth}{0.5pt}\end{center}

\section{Referencing the Blocks}\label{referencing-the-blocks}

You can reference the blocks defined above using their labels:

\begin{itemize}
\tightlist
\item
  The main result is stated in Theorem \hyperref[rthmpythag]{2}.
\item
  Checkout Corollary \hyperref[Rcor-pythag]{3}.\\
\item
  The proof uses Lemma \hyperref[Rlem-euclid]{4}.\\
\item
  Then there's Definition \hyperref[Rdef-orthog]{1}.\\
\item
  Finally, see Video \hyperref[Rvid-yt]{2}.
\end{itemize}

\section{Customization}\label{customization}

Each box has to define four colours. In the file mentioned above the
light-mode border and background are set. For those using dark-mode you
also need to define a dark mode pair. See
\href{themes/dark-styles.scss}{dark-stylesheet} \textbf{maybe move to
own file}.




\end{document}
